\documentclass[11pt]{article}

\usepackage[margin=1in]{geometry}
\usepackage{amsmath}
\usepackage{amssymb}
\usepackage{dsfont}
\usepackage{booktabs}
\usepackage{graphicx}
\usepackage{hyperref}
\usepackage[section]{placeins}
\usepackage[numbers]{natbib}
\usepackage{algorithm}
\usepackage{algorithmic}
\renewcommand{\algorithmicrequire}{\textbf{Input:}}
\renewcommand{\algorithmicensure}{\textbf{Output:}}

\newcommand{\btheta}{\boldsymbol{\theta}}
\newcommand{\bmu}{\boldsymbol{\mu}}
\newcommand{\bSigma}{\boldsymbol{\Sigma}}
\newcommand{\bepsilon}{\boldsymbol{\epsilon}}
\DeclareMathOperator*{\argmax}{arg\,max}

\title{Learning Inventory Control Policies with Evolution Strategies\\[0.35em]
       {\large\normalfont across Lost-Sales, Dual-Sourcing, and Multi-Echelon Problems}}
\author{Nima Manaf \and Willem van Jaarsveld \and Rob Basten\\[0.3em]
        {\small School of Industrial Engineering, Eindhoven University of Technology, The Netherlands}}
\date{\today}

\begin{document}
\maketitle

\begin{abstract}
\noindent
We study whether a single, generic learning recipe --- evolution strategies optimizing compact,
interpretable policies --- can match or beat strong problem-specific heuristics and published deep
reinforcement learning (DRL) across a range of classical inventory control problems. Using
Covariance Matrix Adaptation Evolution Strategies (CMA-ES) we learn the parameters of small policy
function approximators (linear maps, one-hidden-layer networks, and soft decision trees) for four
problems: lost sales, fixed-ordering-cost lost sales, dual sourcing, and divergent multi-echelon
control with special delivery. A recurring design lever is that the \emph{action parameterization is
part of the policy}: rather than fixing an external action grid, each policy carries its own decoder
(a soft-gated direct or ordinal quantity, or soft-tree leaves) and, where a structured heuristic
exists, learns inside that heuristic's coordinate system.

Across a $24$-instance vanilla lost-sales
surface, learned policies are instance-best in $22$ of $24$ cases, and in $47$ of $48$ on the
$48$-instance fixed-cost surface. On the six dual-sourcing instances of \citet{gijsbrechts2022drl}
the learned soft tree matches the capped dual-index heuristic --- itself a $\le0.11\%$ proxy for the
bounded-DP optimum --- on every instance, and thus sits well inside the band that the published A3C
learner ($0.51$--$1.85\%$ gap) does not reach. On both reported multi-echelon settings a
direct-level soft tree improves on the best in-environment constant base-stock by $\approx 14.4\%$.

The policies carry tens to a few hundred parameters, are trained without per-problem hyperparameter
tuning, and the reported single-state actions are validated against an independent Rust rollout. The
results suggest evolution strategies are a practical, portable, and interpretable alternative for
learning inventory policies.
\end{abstract}

\section{Introduction}
\label{sec:introduction}

Inventory management is a fundamental sequential decision problem in operations management: each
period a controller decides how much to order so as to balance the cost of holding stock against the
cost of shortages. Many such problems are naturally formulated as Markov decision processes (MDPs),
but exact dynamic programming is tractable only for small instances because the state space grows
with the lead time and the number of stocking locations. The literature has therefore relied either
on problem-specific \emph{heuristics} that exploit structural properties --- base-stock, $(s,S)$,
$(s,nQ)$, dual-index, and capped dual-index policies --- or, more recently, on \emph{deep
reinforcement learning} (DRL), which represents the policy as a neural network and tunes it from
simulated experience \citep{gijsbrechts2022drl,boute2022deep}. Heuristics require specialized
analysis and simplifying assumptions that limit their reach; DRL is generic but demands substantial
compute and extensive hyperparameter tuning, and its networks --- typically thousands to millions of
parameters --- are hard to interpret \citep{gijsbrechts2022drl,kaynov2024owmr}.

This paper asks a different question. Rather than proposing a new method for one problem, we ask
whether a single, lightweight learning recipe can serve as a \emph{portable} alternative to bespoke
heuristics across structurally different inventory problems, while staying small enough to interpret
and cheap enough to retrain. The recipe is to optimize the parameters of a compact policy function
approximator with Covariance Matrix Adaptation Evolution Strategies (CMA-ES), a gradient-free,
population-based optimizer that has been shown to scale to neural-network policies
\citep{Hansen2001CompletelyStrategies,Salimans2017EvolutionLearning}. CMA-ES is attractive here for
three reasons that hold across all four problems we study. First, it optimizes the
realized objective --- the simulated long-run average cost --- directly, with no value-function
bootstrapping, reward shaping, or credit assignment to individual actions, so the same loop applies
unchanged whether the action is a scalar order, a regular/expedite pair, or a pair of echelon
order-up-to levels. Second, its built-in covariance adaptation supplies exploration automatically,
removing the exploration-strategy design that DRL requires. Third, because it never differentiates
the policy, it accommodates non-smooth, discretized, and clipped action maps that are awkward for
backpropagation but natural for inventory controls.

A second, methodological observation runs through every problem we study: \emph{the action
parameterization is part of the policy, not a fixed property of the environment}. A learner that can
only emit a raw scalar order is implicitly restricted to whatever that scalar can express; a learner
whose decoder is shaped like the relevant structured heuristic --- an ordinal quantity for lost
sales, capped-dual-index coordinates for dual sourcing, direct order-up-to levels for multi-echelon
--- searches a far more useful action geometry. We treat this decoder choice as a first-class design
dimension and show that it, rather than network size, drives several of the results, including the
fixed-cost order/no-order decision and the multi-echelon operating region.

Our contributions are as follows.
\begin{enumerate}
    \item \textbf{A cross-problem evolution-strategies result.} We show that CMA-ES-trained compact
    policies match or beat strong heuristics and published DRL on four classical problems --- lost
    sales, fixed-cost lost sales, dual sourcing, and divergent multi-echelon control --- using one
    optimization loop and per-problem horizons of thousands to tens of thousands of periods, without
    per-problem hyperparameter tuning. Learned policies are instance-best in $22/24$ vanilla and
    $47/48$ fixed-cost lost-sales instances, match the capped dual-index optimal proxy on all six
    dual-sourcing instances, and improve on the best in-environment constant base-stock by
    $\approx 14.4\%$ on both multi-echelon settings.
    \item \textbf{Action parameterization as policy design.} We formalize and evaluate a family of
    policy-owned decoders (soft-gated direct quantity, soft-gated ordinal quantity, soft trees with
    linear leaves) and per-problem action geometries (capped-dual-index coordinates; direct-level
    vs.\ reduced-grid multi-echelon actions), and show that the decoder, not the backbone width, is
    what lets a small policy reach the heuristic's operating region.
    \item \textbf{Compact, interpretable, and verifiable policies.} The learned policies carry tens
    to a few hundred parameters --- orders of magnitude fewer than published DRL networks --- and we
    give a worked example mapping concrete states to orders for each decoder, with the reported
    single-state actions independently validated against a Rust rollout.
\end{enumerate}

The remainder of the paper is organized as follows. Section~\ref{sec:related} reviews related work
on learned inventory policies, evolution strategies, dual sourcing, and multi-echelon control.
Section~\ref{sec:methods} presents the unified methodology: the policy function approximators and
their decoders, the per-problem action geometries, the CMA-ES optimizer and evaluation protocol, and
a worked state-to-order example. Section~\ref{sec:numerical} reports the numerical study across the
four problems, with interpretation of where each decoder wins and where heuristics still lead.
Section~\ref{sec:discussion} draws managerial insights, and Section~\ref{sec:conclusion} concludes
with limitations and future work.

\section{Related work}
\label{sec:related}

Classical inventory problems are easily cast as MDPs but are typically intractable by exact dynamic
programming because of the curse of dimensionality, motivating both analytical heuristics and,
recently, learned policies \citep{Zipkin2008OldSystems,Bijvank2011LostSalesReview,boute2022deep}. A
central idea in the DRL line is to represent the policy and/or value function with a neural network
that maps states to a distribution over order quantities \citep{boute2022deep}. Deep Q-Networks were
applied early to the beer game and to perishable inventory; A3C was used to demonstrate cross-problem
flexibility on dual sourcing, lost sales, and one-warehouse multi-retailer control
\citep{gijsbrechts2022drl}, and proximal policy optimization has since been adopted for joint
ordering, asymmetric one-warehouse multi-retailer systems, and non-stationary demand
\citep{kaynov2024owmr}. These methods share a state--action--reward learning structure: rewards over
a trajectory are attributed to individual actions, and network parameters are tuned to promote
favorable actions. Their generality comes at the cost of heavy tuning and large, opaque networks.

Evolution strategies take a different route. CMA-ES is a gradient-free, population-based optimizer
that adapts a full covariance matrix to search complex, high-dimensional, ill-conditioned landscapes
\citep{Hansen2001CompletelyStrategies}; although not originally a policy-learning method, it was
shown to train neural-network control policies at scale as a competitive alternative to gradient-based
DRL \citep{Salimans2017EvolutionLearning}. Within inventory control, evolutionary methods have
chiefly been used to tune the parameters of fixed-form heuristics rather than to learn generic
policies; learning generic ordering policies with evolution strategies, and doing so across several
structurally distinct problems, is the gap this paper addresses.

The problems we study anchor our comparisons to well-developed literatures. For \emph{lost sales},
myopic and base-stock heuristics and the capped base-stock policy are strong benchmarks
\citep{Zipkin2008OldSystems,Xin2021CappedBaseStock}, and the fixed-ordering-cost variant adds
$(s,S)$, $(s,nQ)$, and modified $(s,S,q)$ comparators \citep{Bijvank2015ParametricPolicies}. For
\emph{dual sourcing}, the structured policies are the single-index, dual-index
\citep{veeraraghavan2008dualindex}, capped dual-index, and tailored base-surge heuristics
\citep{sheopuri2010dualsourcing}; \citet{gijsbrechts2022drl} benchmark A3C against these and against a
bounded-DP optimum on small instances. For divergent \emph{multi-echelon} control with partial lost
sales, the relevant family includes the neuro-dynamic-programming study of \citet{vanroy1997ndp}, the
partial-lost-sales two-echelon model of \citet{nahmias1994partial}, multilocation approximations
\citep{federgruen1984multilocation}, the broad typology of \citet{dekok2018typology}, and recent DRL
for one-warehouse multi-retailer systems \citep{kaynov2024owmr}. We use these as external benchmarks
and otherwise train directly for each problem.

\section{Methods}
\label{sec:methods}

This section states once the modeling principle that runs through all four problems, describes the
policy function approximators and their decoders, gives the per-problem action geometries, and then
the CMA-ES optimizer and the evaluation protocol. The per-problem environment definitions appear with
the experiments in Section~\ref{sec:numerical}.

\subsection{The action parameterization is part of the policy}
\label{sec:methods-principle}

Across every problem we treat the map from latent policy outputs to a feasible order as a learned
component of the policy, not a fixed restriction imposed by the environment. The environment emits a
raw state; the policy applies its own divide-by-scale normalization and then a \emph{decoder} that
turns latent outputs into an integer order (lost sales), a regular/expedite order pair (dual
sourcing), or a pair of echelon order-up-to levels (multi-echelon). Where a structured heuristic
exists, we let the decoder live in that heuristic's coordinate system so the search explores a useful
action geometry rather than an arbitrary raw-quantity space. This single principle replaces the
per-problem boilerplate and makes the decoder a first-class design dimension we evaluate alongside the
backbone.

\subsection{Policy function approximators}
\label{sec:methods-policies}

Every learned policy composes a \emph{backbone} that maps the normalized state feature vector
$x\in\mathbb{R}^d$ to latent outputs with a \emph{decoder} that turns those outputs into the action.
We write $\sigma$ for the logistic function and $\mathrm{softplus}(z)=\ln(1+e^{z})$, and we
denote nearest-nonnegative-integer rounding by
$\lfloor z\rceil_{\ge 0} := \max\{0,\lfloor z\rceil\}$, where $\lfloor\cdot\rceil$ rounds to the
nearest integer. We use two backbones: a linear map $u_\theta(x)=Wx+b$, and a
one-hidden-layer network $u_\theta(x)=V\phi(Ux+c)+e$ of width $H=8$ with the self-normalizing SELU
activation $\phi$ \citep{Klambauer2017Self}. The benchmark surface pairs the backbones with three
policy-owned decoders.

\paragraph{Soft-gated direct quantity (L-SGD, NN-SGD).}
The backbone emits a gate score $g(x)$ and a quantity score $q(x)$, and the order is the gated
softplus quantity
\[
a = \big\lfloor\sigma(g(x))\,\mathrm{softplus}(q(x))\big\rceil_{\ge 0}.
\]
The gate decides \emph{whether} to order and the softplus head decides \emph{how much}, a
factorization that matters most under a fixed ordering cost.

\paragraph{Soft-gated ordinal quantity (L-SGO, NN-SGO).}
The backbone emits a gate score $g(x)$ and $M$ ordinal threshold scores $o_1(x),\dots,o_M(x)$, and
the order gates a sum of threshold probabilities,
\[
a = \Big\lfloor\sigma(g(x))\sum_{k=1}^{M}\sigma(o_k(x))\Big\rceil_{\ge 0},
\]
so the threshold count $M$ sets both the largest representable quantity and the decoder's parameter
count. We use $M=20$ thresholds.

\paragraph{Soft trees with linear leaves (Tree-1, Tree-2).}
A soft decision tree with oblique splits computes leaf path-probabilities $\pi_\ell(x)$; each leaf
emits a local softplus-affine quantity $q_\ell(x)=\mathrm{softplus}(\alpha_\ell^\top x+\beta_\ell)$,
and the order is the rounded mixture
\[
a = \Big\lfloor\sum_\ell \pi_\ell(x)\,q_\ell(x)\Big\rceil_{\ge 0}.
\]
We report depth-1 (Tree-1) and depth-2 (Tree-2) trees. The same soft-tree machinery, with
constant or vector-valued leaves, supplies the structured decoders for dual sourcing and
multi-echelon control below.

\paragraph{Parameter counts.}
For an $L=4$ lost-sales pipeline state ($d=4$, $d+1=5$), each learned policy's trainable parameter
count follows in closed form from its decoder. The linear and network direct-quantity policies carry
$2(d+1)=10$ and $H(d+1)+2(H+1)=58$ parameters; the depth-1 and depth-2 soft trees carry $3(d+1)=15$
and $7(d+1)=35$; and with $M=20$ thresholds the ordinal policies carry $(M+1)(d+1)=105$ (L-SGO) and
$H(d+1)+(M+1)(H+1)=229$ (NN-SGO). These are orders of magnitude below the thousands-to-millions of
parameters typical of DRL policies for the same problems \citep{gijsbrechts2022drl,kaynov2024owmr}.

\subsection{Per-problem action geometries}
\label{sec:methods-geometries}

\paragraph{Lost sales and fixed-cost lost sales.}
The feature vector is the pipeline state of dimension $d=L$; the decoder emits a single integer order
$a\in\mathbb{Z}_{\ge0}$. Under a fixed ordering cost the gated decoders expose the order/no-order
boundary explicitly through their gate.

\paragraph{Dual sourcing.}
The leaves emit a control in \emph{factorized capped-dual-index coordinates}
$(s_e,\ \Delta_r,\ \bar c_r)$ with $s_r=s_e+\Delta_r$, where $s_e$ is the expedite order-up-to level,
$s_r$ the regular order-up-to level, and $\bar c_r$ a cap on the regular order. The expedited order
raises net inventory toward $s_e$ and the regular order raises the inventory position toward $s_r$
subject to the cap. We compare a small-cap variant with $\bar c_r$ on the discrete grid
$\{1,2,3,4,6,8,12\}$ (axis-aligned, constant leaves) and a continuous capped-dual-index-targets
variant. Both keep the policy in the heuristic's coordinate system while letting the targets depend
on state; a raw direct-order output cannot express the capped dual-index control.

\paragraph{Multi-echelon.}
The two leaf outputs map to a state-dependent warehouse order-up-to level $y^w_{S}$ and a shared
retailer order-up-to level $y^r_{S}$. We contrast two action designs. The \textbf{grid} design draws
the levels from the reduced discrete set used by \citet{gijsbrechts2022drl},
$y^w\in\{50,60,\dots,100\}$ and $y^r\in\{0,5,\dots,40\}$; the \textbf{direct} design (ours) lets the
leaves estimate the levels as continuous values rounded to non-negative integers, bounded only by the
physical caps $y^w\in[0,C^w]$, $y^r\in[0,C^r]$. The distinction is consequential because the
cost-minimizing warehouse base-stock for these settings is far above the reduced grid's cap of $100$;
the grid is a policy artifact, not a problem constraint.

\subsection{Optimization with CMA-ES}
\label{sec:methods-cmaes}

We optimize the policy parameter vector $\btheta$ with Covariance Matrix Adaptation Evolution
Strategies, a gradient-free, population-based optimizer
\citep{Hansen2001CompletelyStrategies,Salimans2017EvolutionLearning}. CMA-ES models the population
distribution as a multivariate Gaussian $\mathds{P}_\psi(\btheta)\sim\mathcal{N}(\bmu,\sigma^2\bSigma)
=\bmu+\sigma\,\mathcal{N}(0,\bSigma)$ and, at each iteration, samples a population, evaluates each
member's simulated average cost, and updates $(\bmu,\sigma,\bSigma)$ from the best members
(Algorithm~\ref{alg:cmaes-framework}). The mean is updated by weighted recombination of the
$N'$ best members,
\[
\bmu_{i+1}=\bmu_i+\alpha_\mu\sum_{n'=1}^{N'} w^{n'}\big(\btheta^{n'}_{i+1}-\bmu_i\big),
\qquad \sum_{n'} w^{n'}=1,
\]
the covariance is adapted by combined rank-one (evolution-path) and rank-$\mu$ updates with
exponential smoothing, and the global step size $\sigma$ is adapted online by cumulative step-length
adaptation; we use the reference implementation of \citet{Hansen2019pycma}. Relative to A3C/PPO-style
DRL, CMA-ES has few hyperparameters, supplies its own exploration through $\bSigma$, and optimizes
the realized objective directly, so the loop is unchanged across our four problems.

\begin{algorithm}[!htbp]
\caption{CMA-ES policy optimization (used for all four problems)}
\label{alg:cmaes-framework}
\begin{algorithmic}[1]
    \STATE \textbf{Input:} population size $N$, initial mean $\bmu_1$, step size $\sigma$, simulator and policy decoder
    \FOR{iteration $i=1,2,\ldots$}
        \STATE Sample $\btheta^{(n)}_i\sim\mathcal{N}(\bmu_i,\sigma_i^2\bSigma_i)$ for $n=1,\dots,N$
        \STATE Evaluate average cost $\hat C_T(\btheta^{(n)}_i)$ by rolling out the decoded policy in the simulator
        \STATE Update $\bmu_{i+1},\sigma_{i+1},\bSigma_{i+1}$ from the $N'$ lowest-cost members (weighted recombination; rank-one and rank-$\mu$ covariance updates; cumulative step-length adaptation)
    \ENDFOR
    \STATE \textbf{Output:} best evaluated policy parameters $\btheta^\star$
\end{algorithmic}
\end{algorithm}

\paragraph{Warm-starting (dual sourcing).}
When a strong static control exists in the policy's coordinate system, we warm-start CMA-ES at it. For
dual sourcing the best static capped-dual-index control is near-optimal, so we initialize $\bmu_1$ at
the encoded capped-dual-index solution and let CMA-ES refine around it; without the warm start the
search occasionally lands in a worse basin. The improvements are thus a matter of action geometry and
reliable optimization, not extra capacity.

\paragraph{Divergence handling.}
CMA-ES runs occasionally diverge, producing a degenerate policy with a runaway cost. We detect and
exclude such runs rather than report them; in the result tables a dagger ($\dagger$) marks an excluded
diverged run, and the comparison is made over the remaining policies.

\paragraph{Evaluation protocol.}
Policies are evaluated by simulation under protocols matched to each problem. For \emph{lost sales}
and \emph{fixed-cost lost sales} we train with CMA-ES (population $64$, training horizon $2000$) and
evaluate the long-run average cost over a horizon of $10^6$ periods averaged across $10$ consecutive
seeds. For \emph{dual sourcing} we use a common-random-number (CRN) protocol --- the same $70$
demand-path seeds and horizon $6\times10^4$ for every policy --- so the dominant per-seed variance
cancels and the paired cost difference is statistically meaningful (paired SEM $\sim0.002$--$0.006$
versus per-seed standard deviation $\sim0.5$). For \emph{multi-echelon} we evaluate the long-run
average cost over a horizon of $3\times10^4$ periods and $6$ seeds.

\subsection{From state to order: a worked example}
\label{sec:methods-worked}

To make the decoders concrete we trace how each maps a fixed state to an order, using validated
numbers reproduced by the script \texttt{paper/policy\_state\_action\_examples.py}. For lost sales we
use the canonical instance \texttt{lit\_poisson\_p4\_l4} ($L=4$, $h=1$, $p=4$, Poisson mean $5$). The
policy input is the $4$-vector $[I_{\mathrm{eff}},q_{t-3},q_{t-2},q_{t-1}]$: coordinate $0$ is the
effective inventory at decision time (the environment folds on-hand into the oldest in-transit slot)
and coordinates $1$--$3$ are the three still-outstanding pipeline orders; the policy normalizes by
$\texttt{state\_scale}=20$ before acting and the action is the order quantity. Table~\ref{tab:worked}
shows the order each decoder returns for two states, a near-empty system $[2,0,0,0]$ and a
well-stocked, full-pipeline system $[8,5,5,5]$.

\begin{table}[!htbp]
\centering
\small
\caption{Worked example: order returned by each learned lost-sales decoder for two states of the
\texttt{lit\_poisson\_p4\_l4} instance ($L=4$, $h=1$, $p=4$, Poisson mean $5$). State coordinate $0$
is effective inventory; coordinates $1$--$3$ are the outstanding pipeline. Values are reproduced by
\texttt{paper/policy\_state\_action\_examples.py} and, for the soft trees, equal the exposed Rust
forward by construction.}
\label{tab:worked}
\begin{tabular}{lcc}
\toprule
Decoder & State $[2,0,0,0]$ & State $[8,5,5,5]$ \\
\midrule
Tree-1 (depth-1 linear leaf) & $11$ & $3$ \\
Tree-2 (depth-2 linear leaf) & $14$ & $3$ \\
L-SGD (linear soft-gated direct) & $12$ & $3$ \\
NN-SGD (network soft-gated direct) & $6$ & $4$ \\
L-SGO (linear soft-gated ordinal) & $7$ & $4$ \\
NN-SGO (network soft-gated ordinal) & $6$ & $4$ \\
\bottomrule
\end{tabular}
\end{table}
\FloatBarrier

The arithmetic behind two representative cells illustrates the decoders. For state $[2,0,0,0]$ the
linear soft-gated direct policy (L-SGD) opens its gate, $\sigma(7.435)=0.9994$, and emits
$\mathrm{softplus}(11.888)=11.888$; the product is $11.881$ and the decoder returns
$\lfloor 11.881\rceil_{\ge 0}=12$. The depth-2 soft tree
routes the same state almost entirely to one leaf (path probabilities
$[0,0,1.0,0]$) whose linear-leaf quantity is $14.048$, giving the order
$\lfloor 14.048\rceil_{\ge 0}=14$; this value equals the
exposed Rust forward \texttt{soft\_tree\_action\_from\_flat\_params} at that state by construction.
For the well-stocked state $[8,5,5,5]$ every decoder backs off to a near-replenishment order of $3$
or $4$, as expected when the pipeline already covers near-term demand.

For \emph{dual sourcing} we use \texttt{dual\_l2\_ce105} ($l_r=2$), where the best model is a
depth-2 axis-aligned constant-leaf soft tree with the small-cap capped-dual-index adapter. The
reduced post-decision state is $[I_t, q_{t-1}]$ (net inventory and the regular pipeline). A soft-tree
leaf emits a $3$-control $(s_e,\Delta_r,\bar c_r)$ on the discrete grid, and the adapter maps it to
orders: $s_r=s_e+\max(\Delta_r,0)$, $\text{expedite}=\mathrm{clip}(s_e-I_t,0,12)$, and
$\text{regular}=\min(\max(0,s_r-(I_t+q_{t-1})-\text{expedite}),\bar c_r,12)$. For state $[6,0]$ the tree
selects the leaf control $(4,3,2)$ (dominant leaf probability $0.87$), and the adapter returns
$(\text{expedite}=0,\ \text{regular}=1)$; for the leaner state $[2,0]$ the same leaf control yields
$(\text{expedite}=2,\ \text{regular}=2)$ --- the expedite supplier covers the immediate net-inventory
shortfall while the capped regular order rebuilds the position.

For \emph{multi-echelon} we use \texttt{gijsbrechts2022\_setting1}, where the best run is a depth-2
oblique soft tree with linear leaves in the direct-level action mode. The raw decision state has
dimension $22$ (warehouse on-hand-after-arrival and its pipeline, then per-retailer on-hand and
pipelines for $10$ retailers), normalized by $\texttt{state\_scale}=100$, and the two leaf outputs are
the warehouse and retailer order-up-to levels directly. For a state with the warehouse holding $300$
on hand (empty pipeline) and each of the $10$ retailers holding $20$ on hand with a pipeline of $5$,
the dominant leaf (probability $1.0$) emits continuous levels $[298.572,54.445]$, rounded and clamped
to a warehouse order-up-to of $299$ and a retailer order-up-to of $54$.

\paragraph{Validation.}
All six lost-sales policies and the dual-sourcing policy were validated by a full Python-vs-Rust
rollout cost match on a fixed Poisson$(5)$ (lost sales) / Uniform$\{0,\dots,4\}$ (dual) demand
sequence of $100{,}000$ periods; every absolute cost difference was exactly $0.000$ (well under the
$0.05$ tolerance), confirming the divide-by-scale normalization handling, and no reported action was
withheld. The lost-sales soft-tree single-state actions are the exposed Rust forward by construction,
and the script additionally confirms the Python tree forward agrees with the Rust function at each
printed state. The multi-echelon model has no deterministic demand-binding and a stochastic
warehouse-allocation/emergency-shipment rollout, so a full cost match was not run; instead the vector
soft-tree forward was validated by extracting the warehouse action dimension into a scalar control
and matching it against the exposed Rust scalar forward over $5000$ random raw states (maximum
$|{\rm py}-{\rm rust}|=0$). The retailer dimension uses the identical leaf math and differs only by a
zero min-value offset and an upper clamp, so it is faithful by construction; the single multi-echelon
action is therefore reported with this cross-check rather than a rollout cost match.

\section{Numerical study}
\label{sec:numerical}

We now report the four problems as evidence for the unifying claim laid out in
Section~\ref{sec:introduction}: that one CMA-ES recipe over compact, policy-owned decoders carries
across structurally distinct inventory problems. Each subsection states the problem, then reports
results with interpretation of where a decoder wins and where it does not.

\subsection{Lost sales and fixed-cost lost sales}

\subsubsection{Problem description}

We study single-item lost-sales inventory control and its fixed-ordering-cost variant. Each period
$t$ the controller observes the inventory state --- the on-hand level and the outstanding-order
pipeline $(q_{t-L+1},\dots,q_{t-1})$ --- and chooses an integer replenishment order
$q_t\in\mathbb{Z}_{\ge 0}$ that arrives after a fixed lead time $L$. Demand $D_t$ is stochastic;
sales are limited to on-hand stock and \emph{unmet demand is lost} rather than backordered. With
on-hand stock $x_t$ available to serve demand $D_t$, the period ends with inventory
$I_t=(x_t-D_t)^{+}$ and lost sales $(D_t-x_t)^{+}$ (writing $(\cdot)^{+}=\max(\cdot,0)$), incurring
the one-period cost
\[
c_t = h\,(x_t-D_t)^{+} + p\,(D_t-x_t)^{+},
\]
a holding cost $h$ per unit left on hand and a penalty $p$ per unit of lost sales.

\textbf{Fixed-cost lost sales} is a simple variant of the same problem: the dynamics, state, and
lost-sales rule are identical, and the one-period cost gains a single term
\[
c_t^{K} = c_t + K\,\mathbf{1}\{q_t>0\},
\]
a fixed cost $K$ charged whenever a positive order is placed. This term makes the
\emph{order/no-order} decision --- whether to replenish at all --- the dominant one, and it is
precisely on this decision that the action parameterizations below differ most.

\subsubsection{Heuristic comparators}

The learned columns are the six policies of Section~\ref{sec:methods-policies} (L-SGD, NN-SGD, L-SGO,
NN-SGO, Tree-1, Tree-2) with the lost-sales action geometry of Section~\ref{sec:methods-geometries};
we describe only the problem-specific heuristic comparators here.

The vanilla heuristic columns are Myopic-1 (M1), Myopic-2 (M2), and the standard vector base-stock
policy (SVBS). M1 minimizes the lead-time-deep myopic action value with no future ordering; M2 adds
one discounted continuation step using the next-period M1 action; SVBS computes vector base-stock
levels over the remaining lead-time positions and orders the minimum raise implied by those levels.
The fixed-cost heuristic columns use the inventory position
\[
P_t=x_t+\sum_{j=1}^{L-1} q_{t-j}.
\]
The $(s,S)$ policy orders $S-P_t$ when $P_t\le s$ and otherwise orders zero. The $(s,nQ)$ policy
orders the smallest positive multiple $nQ$ that raises $P_t+nQ$ above $s$, again only when
$P_t\le s$. The modified $(s,S,q)$ policy orders $\min\{q,S-P_t\}$ when $P_t\le s$ and zero
otherwise. Thus the setup-cost comparators expose their order/no-order threshold and their own
batching or capped-raise parameters as part of the policy definition. The learned-policy parameter
counts for the $L=4$ pipeline state are given in Section~\ref{sec:methods-policies}.

\subsubsection{Experiment setup}

\paragraph{Instances.}
We benchmark on a multi-instance surface rather than a single example. The vanilla lost-sales
surface has holding cost $h=1$, mean demand $\mu=5$, lead times $L\in\{4,6,8,10\}$, lost-sales
penalties $p\in\{4,19\}$, and the three demand families below, giving $3\times4\times2=24$
instances. The fixed-cost problem uses the same vanilla surface and adds the two setup costs
$K\in\{5,25\}$, giving $48$ fixed-cost instances.

\paragraph{Demand families.}
The three reported families are mean-preserving, all with $\mathbb{E}[D_t]=5$:
\begin{itemize}
    \item Poisson demand with mean $5$,
    \item Geometric demand with mean $5$,
    \item positively correlated two-state Markov-modulated Poisson (MMPP2$+$) demand with
    $(\lambda_{\mathrm{low}},\lambda_{\mathrm{high}})=(3,7)$ and $p_{00}=p_{11}=0.9$.
\end{itemize}
For the MMPP2$+$ case the latent regime $S_t\in\{\mathrm{low},\mathrm{high}\}$ follows a two-state
Markov chain with stay-probabilities $p_{00},p_{11}$, and $D_t\mid S_t\sim
\mathrm{Poisson}(\lambda_{S_t})$. Because $p_{00}=p_{11}$ the stationary regime probabilities are
$1/2$ each, preserving the mean $\tfrac12(3)+\tfrac12(7)=5$. The temporal dependence is set by the
regime eigenvalue $\rho_{\mathrm{regime}}=p_{00}+p_{11}-1=0.8$; after Poisson observation noise
the implemented lag-one demand autocorrelation is
\[
\mathrm{Corr}(D_t,D_{t+1})
=\frac{\pi_{\mathrm{low}}\pi_{\mathrm{high}}(\lambda_{\mathrm{high}}-\lambda_{\mathrm{low}})^2
(p_{00}+p_{11}-1)}{\mathrm{Var}(D_t)}=\frac{3.2}{9}\approx0.356,
\]
whereas the Poisson and Geometric families are i.i.d.\ (zero autocorrelation at every lag).

\paragraph{Training and evaluation.}
We use the lost-sales protocol of Section~\ref{sec:methods-cmaes}: CMA-ES training (population $64$,
horizon $2000$) and evaluation of the long-run average cost over horizon $10^6$ with $10$ consecutive
seeds. The benchmark heuristics are the myopic-1, myopic-2, and standard vector base-stock (SVBS)
policies for vanilla lost sales, and the $(s,S)$, $(s,nQ)$, and modified $(s,S,q)$ policies for the
fixed-cost variant, using literature-reported parameters where available
\citep{Zipkin2008OldSystems,Bijvank2015ParametricPolicies}.

\subsubsection{Results}

Tables~\ref{tab:results-vanilla-lost-sales} and~\ref{tab:results-fixed-cost-lost-sales} report the
mean cost of the six learned policy approximators against the heuristic baselines across the reported
vanilla and fixed-cost benchmark surfaces. The tables use the same layout: policy rows, lead-time
column groups, and demand-family subcolumns. A dagger marks a diverged CMA-ES run (excluded from the
comparison); bold marks the best reported heuristic or learned policy in each instance column. The
fixed-cost table repeats the vanilla instance surface and adds the two setup costs $K\in\{5,25\}$.

\begin{table}[!htbp]
\centering
\scriptsize
\setlength{\tabcolsep}{3pt}
\caption{Vanilla lost-sales benchmark: average cost by policy, shortage cost, lead time, and demand family across the reported 24-instance surface.}
\label{tab:results-vanilla-lost-sales}
\resizebox{\textwidth}{!}{%
\begin{tabular}{llrrrrrrrrrrrr}
\toprule
 &  & \multicolumn{3}{c}{Lead time $L=4$} & \multicolumn{3}{c}{Lead time $L=6$} & \multicolumn{3}{c}{Lead time $L=8$} & \multicolumn{3}{c}{Lead time $L=10$} \\
\cmidrule(lr){3-5}\cmidrule(lr){6-8}\cmidrule(lr){9-11}\cmidrule(lr){12-14}
Shortage $p$ & Policy & Pois & Geom & MMPP2$+$ & Pois & Geom & MMPP2$+$ & Pois & Geom & MMPP2$+$ & Pois & Geom & MMPP2$+$ \\
\midrule
$4$ & Optimal & 4.730 & 10.610 &  &  &  &  &  &  &  &  &  &  \\
 & M1 & 5.065 & 11.217 & 9.823 & 5.422 & 11.650 & 11.189 & 5.698 & 11.924 & 12.256 & 5.890 & 12.075 & 13.136 \\
 & M2 & 4.819 & 10.800 & 9.452 & 5.051 & 11.080 & 10.591 & 5.199 & 11.270 & 11.460 & 5.310 & 11.400 & 12.155 \\
 & SVBS & 5.836 & 13.728 & 10.312 & 6.781 & 15.738 & 12.271 & 7.694 & 17.632 & 13.995 & 8.038 & 19.421 & 15.190 \\
 & L-SGD & 4.750 & 10.645 & 8.571 & 4.897 & 10.790 & 9.171 & \textbf{4.974} & \textbf{10.850} & 9.445 & 5.048 & 10.921 & 9.563 \\
 & NN-SGD & 4.774 & 10.687 & 8.565 & 4.900 & 10.791 & 9.643 & 5.033 & 10.983 & 9.674 & 5.066 & 10.993 & 9.697 \\
 & L-SGO & 4.768 & 10.647 & \textbf{8.549} & \textbf{4.893} & 10.823 & \textbf{9.109} & $\dagger$ & 10.884 & \textbf{9.443} & 16.978 & 10.991 & 9.725 \\
 & NN-SGO & 4.786 & 10.692 & 8.577 & 5.064 & 10.799 & 9.183 & 4.978 & $\dagger$ & 9.481 & $\dagger$ & 10.995 & 9.726 \\
 & Tree-1 & \textbf{4.749} & \textbf{10.621} & 8.579 & 4.895 & \textbf{10.776} & 9.159 & 4.975 & 10.857 & 9.445 & \textbf{5.013} & 10.934 & 9.590 \\
 & Tree-2 & 4.750 & 10.670 & 8.595 & 4.936 & 10.820 & 9.302 & 4.977 & 10.892 & 9.446 & 5.067 & \textbf{10.912} & \textbf{9.547} \\
\addlinespace
$19$ & Optimal & 8.890 & 22.950 &  &  &  &  &  &  &  &  &  &  \\
 & M1 & 9.173 & 23.836 & 17.182 & 10.269 & 25.821 & 20.129 & 11.092 & 27.407 & 22.399 & 11.820 & 28.672 & 24.218 \\
 & M2 & 8.953 & 23.280 & 17.644 & 9.831 & 24.930 & 20.474 & 10.511 & 26.210 & 22.552 & 11.090 & 27.270 & \textbf{23.913} \\
 & SVBS & 9.565 & 26.429 & 16.531 & 11.234 & 30.103 & 19.677 & 12.333 & 33.008 & \textbf{22.173} & 13.336 & 35.761 & 24.457 \\
 & L-SGD & 8.932 & \textbf{22.993} & 16.686 & 9.759 & 24.486 & 19.862 & 10.261 & 25.374 & 22.378 & 10.634 & 26.287 & 24.197 \\
 & NN-SGD & \textbf{8.927} & 23.455 & \textbf{16.341} & 10.131 & 24.483 & 19.650 & 10.430 & 25.757 & 22.211 & 10.772 & 26.129 & 24.406 \\
 & L-SGO & 9.003 & 23.240 & 16.491 & 9.794 & 24.370 & \textbf{19.549} & 10.330 & 25.766 & 22.233 & $\dagger$ & 26.171 & 24.345 \\
 & NN-SGO & 8.974 & 23.140 & 16.435 & 9.758 & 24.474 & 19.693 & 10.490 & 25.662 & 22.303 & 10.733 & \textbf{26.037} & 24.357 \\
 & Tree-1 & 8.936 & 23.011 & 16.697 & \textbf{9.711} & \textbf{24.349} & 19.915 & \textbf{10.259} & \textbf{25.209} & 22.339 & \textbf{10.629} & 26.064 & 24.101 \\
 & Tree-2 & 8.941 & 23.043 & 16.683 & 9.720 & 24.430 & 19.864 & 10.272 & 25.382 & 22.301 & 10.698 & 26.174 & 24.193 \\
\bottomrule
\end{tabular}}
\par\medskip\footnotesize Optimal entries are shown only for verified literature optimum anchors; blanks mean no true optimum anchor is available for that expanded instance. $\dagger$: diverged CMA-ES run (excluded). Bold: best reported heuristic or learned policy in the instance column.
\end{table}
\FloatBarrier

\begin{table}[!htbp]
\centering
\scriptsize
\setlength{\tabcolsep}{3pt}
\caption{Fixed-order-cost lost-sales benchmark: average cost by policy, setup cost, shortage cost, lead time, and demand family across the reported 48-instance surface.}
\label{tab:results-fixed-cost-lost-sales}
\resizebox{\textwidth}{!}{%
\begin{tabular}{lllrrrrrrrrrrrr}
\toprule
 &  &  & \multicolumn{3}{c}{Lead time $L=4$} & \multicolumn{3}{c}{Lead time $L=6$} & \multicolumn{3}{c}{Lead time $L=8$} & \multicolumn{3}{c}{Lead time $L=10$} \\
\cmidrule(lr){4-6}\cmidrule(lr){7-9}\cmidrule(lr){10-12}\cmidrule(lr){13-15}
Setup $K$ & Shortage $p$ & Policy & Pois & Geom & MMPP2$+$ & Pois & Geom & MMPP2$+$ & Pois & Geom & MMPP2$+$ & Pois & Geom & MMPP2$+$ \\
\midrule
$5$ & $4$ & $(s,S)$ & 9.369 & 14.130 & 12.158 & 9.772 & 14.697 &  & 9.965 & 14.839 &  & 10.015 & 15.024 &  \\
 &  & $(s,nQ)$ & 9.211 & 13.968 & 12.113 & 9.385 & 14.177 &  & 9.569 & 14.389 &  & 9.600 & 14.542 &  \\
 &  & $(s,S,q)$ & 9.202 & 13.939 & 12.077 & 9.380 & 14.299 &  & 9.562 & 14.426 &  & 9.826 & 14.524 &  \\
 &  & L-SGD & 8.784 & 14.069 & 12.314 & 9.558 & 13.933 & 12.681 & 9.588 & 14.136 & \textbf{12.722} & 8.934 & 14.023 & 12.818 \\
 &  & NN-SGD & \textbf{8.733} & 13.649 & 13.655 & \textbf{8.871} & \textbf{13.825} & 12.854 & $\dagger$ & 13.910 & 12.940 & 8.933 & \textbf{13.861} & 12.816 \\
 &  & L-SGO & 8.784 & 13.755 & \textbf{11.807} & 9.009 & 13.876 & 12.613 & 8.993 & \textbf{13.847} & 12.751 & 8.943 & 13.942 & 12.808 \\
 &  & NN-SGO & 8.748 & \textbf{13.637} & 11.831 & 8.871 & 13.841 & \textbf{12.363} & 8.909 & 14.144 & 13.010 & 9.153 & 13.866 & 12.817 \\
 &  & Tree-1 & 8.771 & 14.057 & 12.146 & 8.872 & 13.933 & 12.691 & \textbf{8.905} & 14.134 & 12.770 & \textbf{8.931} & 13.871 & \textbf{12.791} \\
 &  & Tree-2 & 8.783 & 13.751 & 11.982 & 9.560 & 13.837 & 12.665 & 8.954 & 13.865 & 12.735 & 8.934 & 14.019 & 12.899 \\
\addlinespace
$5$ & $19$ & $(s,S)$ & 13.889 & 26.888 &  & 14.841 & 28.669 &  & 15.816 & 31.221 &  & 21.667 & 35.693 &  \\
 &  & $(s,nQ)$ & 13.657 & 26.878 &  & 14.694 & 28.872 &  & 15.301 & 29.870 &  & 19.152 & 33.406 &  \\
 &  & $(s,S,q)$ & 13.560 & 26.835 &  & 14.476 & 28.543 &  & 15.618 & 31.067 &  & 21.529 & 35.603 &  \\
 &  & L-SGD & 13.619 & 27.260 & 21.350 & 14.265 & 28.195 & 24.567 & 14.704 & 28.977 & 27.001 & 15.151 & 29.576 & 28.632 \\
 &  & NN-SGD & 13.590 & \textbf{26.531} & \textbf{20.079} & $\dagger$ & 27.941 & 23.306 & 14.770 & $\dagger$ & 25.924 & \textbf{15.030} & 29.822 & \textbf{28.042} \\
 &  & L-SGO & 13.574 & 26.696 & 20.172 & 14.417 & 28.243 & \textbf{23.255} & $\dagger$ & 29.046 & 26.071 & 15.780 & 30.219 & 28.241 \\
 &  & NN-SGO & \textbf{13.475} & 26.602 & 20.208 & \textbf{14.167} & \textbf{27.859} & 23.542 & 14.788 & \textbf{28.936} & \textbf{25.805} & $\dagger$ & 29.879 & 28.079 \\
 &  & Tree-1 & 13.628 & 27.071 & 21.309 & 14.268 & 28.250 & 24.552 & 14.701 & 28.957 & 27.087 & 15.108 & \textbf{29.517} & 28.673 \\
 &  & Tree-2 & 13.613 & 27.078 & 21.325 & 14.260 & 28.177 & 24.642 & \textbf{14.700} & 28.962 & $\dagger$ & 15.044 & 29.615 & 28.847 \\
\addlinespace
$25$ & $4$ & $(s,S)$ & 16.002 & 18.428 & 17.342 & 16.136 & 18.736 &  & 16.331 & 18.925 &  & 16.562 & 19.036 &  \\
 &  & $(s,nQ)$ & 15.840 & \textbf{18.231} & 17.351 & 15.979 & 18.579 &  & 15.837 & 18.649 &  & 16.020 & 18.872 &  \\
 &  & $(s,S,q)$ & 15.877 & \textbf{18.231} & 17.318 & 16.060 & 18.579 &  & 15.806 & 18.648 &  & 16.331 & 18.890 &  \\
 &  & L-SGD & 16.645 & 19.992 & 19.995 & \textbf{15.812} & 19.992 & 19.995 & 16.988 & 19.992 & 19.995 & 17.491 & 19.992 & 18.043 \\
 &  & NN-SGD & \textbf{15.572} & 19.992 & 19.995 & 17.075 & 19.992 & 19.995 & 17.018 & 19.992 & 19.995 & 17.451 & 19.992 & 17.946 \\
 &  & L-SGO & 17.983 & 19.992 & 19.995 & 15.847 & 19.992 & 19.995 & 15.810 & 18.509 & 19.995 & \textbf{15.651} & 19.992 & \textbf{17.786} \\
 &  & NN-SGO & 15.627 & 19.992 & \textbf{17.279} & $\dagger$ & \textbf{18.424} & \textbf{17.840} & \textbf{15.656} & \textbf{18.431} & \textbf{17.753} & 16.271 & \textbf{18.633} & 17.831 \\
 &  & Tree-1 & 20.004 & 19.992 & 19.995 & 18.321 & 19.992 & 19.995 & 19.497 & 19.992 & 19.995 & 17.864 & 19.992 & 19.995 \\
 &  & Tree-2 & 16.203 & 19.992 & 19.995 & 18.167 & 19.992 & 19.995 & 16.632 & 19.992 & 19.995 & 19.710 & 19.992 & 19.995 \\
\addlinespace
$25$ & $19$ & $(s,S)$ & 21.104 & 32.726 &  & 22.073 & 35.015 &  & 25.969 & 38.859 &  & 33.668 & 43.651 &  \\
 &  & $(s,nQ)$ & 21.016 & 32.946 &  & 22.283 & 34.586 &  & 22.664 & 35.812 &  & 25.247 & 38.055 &  \\
 &  & $(s,S,q)$ & 21.016 & 32.772 &  & 22.373 & 35.116 &  & 25.971 & 38.865 &  & 33.995 & 43.652 &  \\
 &  & L-SGD & 21.166 & 32.668 & $\dagger$ & 22.878 & 34.569 & 33.398 & 24.641 & 35.287 & 35.916 & 23.727 & 36.354 & 36.525 \\
 &  & NN-SGD & 20.845 & 32.645 & 27.434 & 21.709 & \textbf{34.021} & \textbf{30.199} & \textbf{22.345} & 35.389 & 32.862 & \textbf{22.772} & $\dagger$ & 34.924 \\
 &  & L-SGO & 21.672 & 32.937 & \textbf{27.386} & 22.447 & 34.256 & 31.004 & 24.816 & 36.025 & 33.133 & 25.213 & 36.400 & 36.811 \\
 &  & NN-SGO & 20.854 & \textbf{32.435} & 27.423 & \textbf{21.598} & 34.185 & 30.495 & 22.752 & \textbf{35.070} & \textbf{32.537} & 23.061 & \textbf{36.005} & \textbf{34.373} \\
 &  & Tree-1 & 21.173 & 32.640 & 28.727 & 21.914 & 35.108 & 33.330 & 22.654 & 35.884 & 35.822 & 23.750 & 36.241 & 37.530 \\
 &  & Tree-2 & \textbf{20.770} & 33.081 & 28.216 & 22.243 & 34.425 & 33.603 & 22.622 & 35.164 & 32.999 & 22.977 & 36.158 & 36.324 \\
\bottomrule
\end{tabular}}
\par\medskip\footnotesize Blank heuristic cells indicate no literature/catalog heuristic value for that MMPP2$+$ instance. $\dagger$: diverged CMA-ES run (excluded). Bold: best reported heuristic or learned policy in the instance column.
\end{table}
\FloatBarrier

The column winners make the comparison easy to read. On the vanilla surface, learned policies are
best in $22$ of $24$ instances; Tree-1 is the most frequent winner ($9$ instances), followed by L-SGO
($5$ instances). The two instances still won by classical baselines are the high-penalty MMPP2$+$
cases with $L=8,10$. On the fixed-cost surface, learned policies are instance-best among the
available policies in $47$ of $48$ instances. NN-SGO is the most frequent winner ($20$ instances),
followed by NN-SGD ($13$ instances); the only heuristic instance win is Geometric $L=4,p=4,K=25$,
where $(s,nQ)$ and modified $(s,S,q)$ tie. Fixed-cost
lost sales is therefore the sharper architecture test: the fixed charge makes the order/no-order
boundary central, and the best-performing decoder changes with demand family and setup cost.

\paragraph{Interpretation.}
The pattern of decoder wins is informative. On the vanilla surface, where the action is a plain
quantity, the soft trees (Tree-1, Tree-2) and the linear ordinal decoder (L-SGO) dominate because a
mild state-dependent raise is all that is needed and the smaller decoders optimize more reliably; the
extra capacity of the network decoders does not pay for itself. The fixed-cost surface inverts this:
once a setup charge $K$ makes the order/no-order decision dominant, the gated decoders --- which
factor the decision into an explicit gate (\emph{whether} to order) and a quantity head
(\emph{how much}) --- win most instances, with the ordinal head (NN-SGO, $20$ wins) edging the
direct head (NN-SGD, $13$ wins) where the quantity must batch into a few discrete levels. Two
limitations are visible directly in the tables. First, at the high setup cost $K=25,\,p=4$ several
learned policies collapse onto a degenerate \emph{never-order} policy: the repeated entries
$19.992$ (Geometric) and $19.995$ (MMPP2$+$) are the cost of a policy that effectively stops ordering
to avoid the setup charge, which Tree-1 and Tree-2 reach across most $K=25,p=4$ instances; the gated
ordinal decoder (NN-SGO) is the one architecture that consistently escapes this collapse and remains
instance-best, underscoring that the decoder, not the optimizer, governs the order/no-order boundary.
Second, the two vanilla instances still won by a classical baseline are the high-penalty
($p=19$) MMPP2$+$ cases at the longest lead times $L=8$ and $L=10$, where SVBS ($L=8$) and Myopic-2
($L=10$) edge the learned policies; here the combination of a regime-switching, autocorrelated demand
and a deep pipeline is the hardest setting for a stationary compact policy, and closing this gap is
an open item we return to in Section~\ref{sec:discussion}.

\subsection{Dual sourcing}
\label{sec:dual-sourcing}

\subsubsection{Problem description}

We study the dual-sourcing inventory model of \citet{gijsbrechts2022drl}: a single item is
replenished from a slow \emph{regular} supplier with lead time $l_r$ and a fast \emph{expedited}
supplier with lead time $l_e=0$, against i.i.d.\ demand. Each period the controller chooses a
regular order $q^r_t$ (unit cost $c_r$) and an expedited order $q^e_t$ (unit cost $c_e>c_r$);
end-of-period net inventory incurs a holding cost $h$ per unit on hand or a backorder penalty $b$
per unit short. Because only the expedite lead time is zero, the relevant post-decision state
reduces to the net inventory plus the outstanding regular-order pipeline
$(I_t, q^r_{t-l_r+1},\dots,q^r_{t-1})$. The classical structured policies are the single-index,
dual-index \citep{veeraraghavan2008dualindex}, capped dual-index, and tailored base-surge
heuristics \citep{sheopuri2010dualsourcing}.

\subsubsection{Policy and comparator}

Dual sourcing tests whether the same recipe carries over when the action is a structured order pair
rather than a scalar. The learned policy is the soft-tree backbone of
Section~\ref{sec:methods-policies} on the reduced pipeline state, with leaves that emit a control in
the factorized capped-dual-index coordinates $(s_e,\Delta_r,\bar c_r)$ of
Section~\ref{sec:methods-geometries}; we use its small-cap (axis-aligned, constant-leaf) and
continuous-target variants and warm-start CMA-ES at the encoded capped-dual-index solution as
described in Section~\ref{sec:methods-cmaes}. The comparison is therefore one of action geometry and
reliable optimization rather than tree capacity, and is made against the classical structured
heuristics introduced above.

\subsubsection{Experiment setup}

We use the six small-scale benchmark instances of \citet{gijsbrechts2022drl} (their Figure~9),
with $l_r\in\{2,3,4\}$ and $c_e\in\{105,110\}$ and the common parameters of
Table~\ref{tab:ds-params}. The published comparators are the optimality gaps, relative to the
bounded-DP optimum, of the single-index, dual-index, capped dual-index, and tailored base-surge
heuristics together with the A3C learner of \citet{gijsbrechts2022drl}: the capped dual-index gap
is $\le0.11\%$ on every row, making it both the best heuristic and an effective optimal proxy,
whereas A3C ranges from $0.51\%$ to $1.85\%$. We therefore benchmark the learned policy against the
\emph{capped dual-index} cost as the optimal proxy rather than our own bounded DP
(see Appendix~\ref{app:ds-validation}). Policies are trained with warm-started CMA-ES and evaluated
under a common-random-number protocol --- the same $70$ demand-path seeds and horizon $6\times10^4$
for all policies --- so that the dominant per-seed variance cancels and the paired cost difference
is statistically meaningful (paired SEM $\sim0.002$--$0.006$ versus per-seed standard deviation
$\sim0.5$).

\begin{table}[!htbp]
\centering
\caption{Dual-sourcing benchmark instances \citep{gijsbrechts2022drl}. Common to all:
$l_e=0$, $c_r=100$, $h=5$, $b=495$, demand $U\{0,\dots,4\}$, order caps $12$.}
\label{tab:ds-params}
\begin{tabular}{lcc}
\toprule
Instances & $l_r$ & $c_e$ \\
\midrule
dual\_l$\{2,3,4\}$\_ce105 & $2,3,4$ & $105$ \\
dual\_l$\{2,3,4\}$\_ce110 & $2,3,4$ & $110$ \\
\bottomrule
\end{tabular}
\end{table}
\FloatBarrier

\subsubsection{Results}

Table~\ref{tab:ds-results} compares the best learned policy against the capped dual-index (CDI)
optimal proxy under the common-random-number protocol. The headline reading is that the learned soft
tree \emph{matches} CDI: it lands within the discrete-grid rounding floor of CDI on four instances
($\le+0.003\%$) and is, by the paired CRN comparison, statistically but negligibly below it on two
($l_r=2,c_e=110$ and $l_r=4,c_e=110$, by $0.009\%$ and $0.041\%$). These tiny negative margins sit
well inside CDI's own published optimality band ($\le0.11\%$ gap to the bounded-DP optimum), so we do
not read them as the learned policy genuinely outperforming the heuristic; the CDI reference numbers
are easy to reproduce, and several controls in the same coordinate system attain similar costs.
\paragraph{Interpretation.}
The salient result is therefore reliability of reproduction rather than a new champion: by living in
the capped-dual-index coordinate system and warm-starting at the CDI solution, CMA-ES recovers the
near-optimal control on every instance. Because CDI is the best static capped-dual-index comparator
and a $\le0.11\%$ proxy for the bounded-DP optimum, the learned policy thereby sits at the same
near-optimal reference level on every instance column --- comfortably inside the band that the
published A3C learner, with optimality gaps of $0.51$--$1.85\%$, does not reach. The action geometry,
not the learner's capacity, is what makes this attainable: a raw direct-order decoder cannot express
the capped dual-index control and does not reach this level.

\begin{table}[!htbp]
\centering
\small
\caption{Dual-sourcing: best learned soft-tree policy vs.\ the capped dual-index (CDI) optimal proxy, common-random-number evaluation ($70$ seeds, horizon $6\times10^4$). $\Delta\%$ is the learned policy's cost relative to CDI (negative is better); the A3C row reports the published optimality gap from \citet{gijsbrechts2022drl}.}
\label{tab:ds-results}
\resizebox{\textwidth}{!}{%
\begin{tabular}{lcccccc}
\toprule
 & \multicolumn{2}{c}{Regular lead time $l_r=2$} & \multicolumn{2}{c}{Regular lead time $l_r=3$} & \multicolumn{2}{c}{Regular lead time $l_r=4$} \\
\cmidrule(lr){2-3}\cmidrule(lr){4-5}\cmidrule(lr){6-7}
Policy / metric & $c_e=105$ & $c_e=110$ & $c_e=105$ & $c_e=110$ & $c_e=105$ & $c_e=110$ \\
\midrule
CDI proxy & 216.806 & 219.821 & 216.900 & 220.483 & 216.914 & 220.879 \\
Learned soft tree & 216.806 & 219.800 & 216.906 & 220.483 & 216.918 & 220.789 \\
Learned vs.\ CDI (\%) & $+0.000$ (match) & $-0.009$ (match) & $+0.003$ (match) & $+0.000$ (match) & $+0.002$ (match) & $-0.041$ (match) \\
Published A3C gap & 0.52\% & 0.80\% & 0.82\% & 0.51\% & 1.85\% & 1.33\% \\
\bottomrule
\end{tabular}}
\par\medskip\footnotesize All six instances match CDI to within its own published optimality band ($\le0.11\%$): four are at or within the discrete-grid rounding floor ($\le+0.003\%$) and two are negligibly below CDI ($-0.009\%$, $-0.041\%$) under the paired CRN comparison. We report these as matches, not improvements.
\end{table}
\FloatBarrier

\subsection{Multi-echelon inventory control: one warehouse, $R$ retailers with special delivery}

\subsubsection{Problem description}

We study the divergent two-echelon system of \citet{vanroy1997ndp}, also used as a deep-RL
benchmark by \citet{gijsbrechts2022drl}: a single capacitated warehouse replenishes $R$ identical
capacitated retailers. The warehouse orders from an external manufacturer with lead time $l_w$ and
ships to retailers with lead time $l_r$. Each period $t$: (i) the warehouse places an order; (ii)
each retailer raises its inventory position toward a target, shipped from current warehouse stock
(allocated to minimize the spread of retailer positions when warehouse stock is short); (iii)
retailer demand $d_i$ is realized and served from retailer on-hand; (iv) each unmet unit requests a
same-day \emph{special delivery} from leftover warehouse stock with probability $P_w$ and is
otherwise lost --- a hybrid of backlogging and lost sales; (v) pipelines advance. The state is the
warehouse and per-retailer on-hand levels and outstanding-order pipelines. With $I^w_t,I^i_t$ the
end-of-period on-hand inventories and $E_t$ the number of expedited units, the period cost is
\[
c_t = h^w I^w_t + h^r \sum_{i=1}^{R} I^i_t + c^w E_t + p\,(\text{lost units}),
\]
with warehouse/retailer holding costs $h^w,h^r$, unit expedite cost $c^w$, and lost-sale penalty
$p$. Capacities bound the per-period production ($C^m$) and the warehouse and retail inventory
positions ($C^w,C^r$). This is a member of the broader divergent / partial-lost-sales multi-echelon
family \citep{nahmias1994partial, federgruen1984multilocation, dekok2018typology, kaynov2024owmr}.

\subsubsection{Policy and action designs}

This is the most demanding test of the action-parameterization principle, because the wrong action
geometry is fatal. The policy is the soft decision tree with oblique splits and linear leaves of
Section~\ref{sec:methods-policies}, consuming the raw decision state with a policy-owned divide-by-scale
normalization; its two leaf outputs are the state-dependent warehouse and retailer order-up-to levels
$y^w_{S},y^r_{S}$, with the warehouse order $q^w = \big[\min(y^w_{S},C^w)-\mathrm{IP}^w\big]^+$ capped
by $C^m$ (where $\mathrm{IP}^w$ is the warehouse inventory position and $[\cdot]^+=\max(\cdot,0)$) and
each retailer order raising its position toward $\min(y^r_{S},C^r)$. We contrast the
reduced \textbf{grid} action design of \citet{gijsbrechts2022drl} with our \textbf{direct} continuous
design as in Section~\ref{sec:methods-geometries}; the grid caps the warehouse level at $100$ while
the cost-minimizing warehouse base-stock here is $\approx 300$--$525$ (the published constant
base-stock uses $y^w=330,460$), so the grid cannot express the operating region and is a policy
artifact, not a problem constraint. We train for the problem itself and use the published numbers only
as a benchmark, independently of the learner in \citet{gijsbrechts2022drl}.

\subsubsection{Experiment setup}

We evaluate on \emph{both} multi-echelon settings of \citet{gijsbrechts2022drl} --- the two complex
case studies of \citet{vanroy1997ndp} --- which are the only multi-echelon instances they report, so
our coverage of their multi-echelon benchmark is complete. Table~\ref{tab:me-params} reproduces
their Table~3 parameters exactly for both settings under our faithful \texttt{gijs\_2022} dynamics;
retailer demand is $\mathcal{N}(\mu,\sigma^2)$ rounded to a non-negative integer.

\begin{table}[!htbp]
\centering
\caption{Multi-echelon experiment settings, reproducing Table~3 of \citet{gijsbrechts2022drl} (the
two complex case studies of \citet{vanroy1997ndp}) parameter-for-parameter under our faithful
\texttt{gijs\_2022} dynamics. These are the only two multi-echelon settings reported by
\citet{gijsbrechts2022drl}; we evaluate both. Lead times $l_w,l_r$; demand mean $\mu$ and standard
deviation $\sigma$; number of retailers $R$; warehouse and retailer holding costs $h^w,h^r$; unit
expedite cost $c^w$; lost-sales penalty $p$; special-delivery success probability $P_w$; production,
warehouse, and retailer capacities $C^m,C^w,C^r$.}
\label{tab:me-params}
\begin{tabular}{lcc}
\toprule
Parameter & Setting 1 & Setting 2 \\
\midrule
\addlinespace[2pt]
Warehouse lead time $l_w$ & $2$ & $5$ \\
Retailer lead time $l_r$ & $2$ & $3$ \\
Demand mean $\mu$ & $5$ & $0$ \\
Demand standard deviation $\sigma$ & $14$ & $20$ \\
Retailers $R$ & $10$ & $10$ \\
\addlinespace[2pt]
Warehouse holding cost $h^w$ & $3$ & $3$ \\
Retailer holding cost $h^r$ & $3$ & $3$ \\
Unit expedite cost $c^w$ & $0$ & $0$ \\
Lost-sales penalty $p$ & $60$ & $60$ \\
Special-delivery success probability $P_w$ & $0.8$ & $0.8$ \\
\addlinespace[2pt]
Production capacity $C^m$ & $100$ & $100$ \\
Warehouse capacity $C^w$ & $1000$ & $1000$ \\
Retailer capacity $C^r$ & $100$ & $100$ \\
\bottomrule
\end{tabular}
\end{table}
\FloatBarrier

The published benchmarks we compare against are the constant base-stock costs $(330,23)\to 1302$
and $(460,22)\to 1449$, the best neuro-dynamic-programming costs $1179$ and $1318$, and the A3C
improvements over constant base-stock of $8.95\%\pm0.13\%$ and $12.09\%\pm0.39\%$
\citep{vanroy1997ndp, gijsbrechts2022drl}. As the in-environment benchmark we compute the best
\emph{constant} (state-independent) base-stock by grid search over the physical operating region
(warehouse up to $\approx C^w$, retailer up to $C^r$) rather than over the reduced grid. Policies
are trained with CMA-ES; we sweep the action design against tree depth $\in\{2,3\}$ and evaluate the
long-run average cost over a horizon of $3\times 10^4$ periods and $6$ seeds.

\subsubsection{Results}

Table~\ref{tab:me-results} reports the best in-environment constant base-stock, the best learned
policy under each action design, and the learned policy's improvement over the constant base-stock,
alongside the published A3C improvement. The direct-estimation policy improves on the best constant
base-stock by $\approx 14.4\%$ on both settings. The grid-action policy, restricted to $y^w\le 100$,
stays about $230\%$ \emph{above} the benchmark.
\paragraph{Interpretation.}
The gap between the two action designs is the central finding: the grid policy fails not because of
the environment or the learner but because its action space cannot reach the operating region. This
is the action-space trap in its starkest form --- the same soft tree, optimizer, and horizon, with
only the decoder's reachable level set changed, swing from $\approx14\%$ better than the constant
base-stock to more than $200\%$ worse. The best tree depth is instance-dependent (depth~2 for
setting~1, depth~3 for setting~2), and the design-by-depth sweep selects it automatically. One
comparison caveat is important: our $14.4\%$ improvement and the published A3C improvements
($8.95\%$, $12.09\%$) are computed against \emph{different} baselines and under \emph{different} cost
conventions --- our improvement is relative to the best in-environment constant base-stock under the
\texttt{gijs\_2022} dynamics, whereas the A3C figures are relative to the published van-Roy constant
base-stock under the original cost convention --- so the cross-method comparison is indicative of the
direct design's strength rather than a strictly like-for-like ranking.

\begin{table}[!htbp]
\centering
\caption{Multi-echelon learned-policy results (\texttt{gijs\_2022} dynamics; mean cost, lower is
better). ``Improvement'' is the direct learned policy's cost reduction relative to the best
constant base-stock; the last row is the published A3C improvement over constant base-stock from
\citet{gijsbrechts2022drl}.}
\label{tab:me-results}
\begin{tabular}{lcc}
\toprule
Policy / metric & Setting 1 & Setting 2 \\
\midrule
Best const.\ base-stock $(y^w,y^r)$ & $(300,25)\to 911.4$ & $(525,25)\to 1137.8$ \\
Grid tree & $3085.7$ & $3795.1$ \\
Direct tree (ours) & $\mathbf{779.8}$ & $\mathbf{973.6}$ \\
Improvement vs.\ best const. & $\mathbf{14.4\%}$ & $\mathbf{14.4\%}$ \\
Published A3C improvement & $8.95\%$ & $12.09\%$ \\
\bottomrule
\end{tabular}
\end{table}
\FloatBarrier

\section{Discussion and managerial insights}
\label{sec:discussion}

Synthesizing across the four problems, three themes emerge. First, a \emph{single, generic learning
recipe} --- CMA-ES over a compact policy --- is portable: the same optimization loop, unchanged,
matches or beats strong heuristics and published DRL on a scalar lost-sales order, a setup-cost
order/no-order decision, a structured dual-sourcing order pair, and a pair of echelon order-up-to
levels. For a practitioner this means one method to maintain rather than a bespoke heuristic per
problem, and retraining when costs or demand change is a matter of rerunning a short simulation, not
re-deriving structure.

Second, the \emph{action parameterization is where most of the leverage lives}. The backbone width
matters far less than the decoder: on vanilla lost sales a small tree or linear ordinal head suffices;
under a fixed cost the gated decoders that separate \emph{whether} from \emph{how much} are what
avoid the never-order collapse; in dual sourcing the capped-dual-index coordinates are what let the
policy reach the near-optimal region a raw-quantity decoder cannot express; and in multi-echelon the
direct-level design is the difference between a $14\%$ improvement and a policy more than $200\%$ worse.
The managerial reading is that choosing the right control representation --- ideally the coordinate
system of the best known heuristic --- is a more reliable lever than enlarging the network.

Third, the learned policies are \emph{small, interpretable, and cheap}. They carry tens to a few
hundred parameters, train without per-problem hyperparameter tuning, and reduce at run time to a small
matrix multiply or a shallow tree traversal, so they can be deployed in standard ERP or spreadsheet
tooling and audited by non-specialists. Where a structured heuristic underlies the decoder (dual
sourcing, multi-echelon), the learned policy's outputs are directly readable as order-up-to levels and
caps. The worked example of Section~\ref{sec:methods-worked} and its Rust cross-check show the
mapping from state to order is concrete and verifiable.

\section{Conclusion, limitations, and future work}
\label{sec:conclusion}

We have shown that evolution strategies --- specifically CMA-ES over compact, policy-owned decoders
--- learn inventory control policies that match or beat strong heuristics and published DRL across
four structurally distinct problems: lost sales, fixed-cost lost sales, dual sourcing, and divergent
multi-echelon control. Relative to DRL, CMA-ES needs few hyperparameters, supplies its own
exploration, and optimizes the realized cost directly, so the same loop transfers across problems; and
the resulting policies are orders of magnitude smaller than published DRL networks while remaining
interpretable and verifiable.

Several limitations bound these claims and point to future work. The comparisons are indicative rather
than strictly like-for-like in two places we flag in the text: the dual-sourcing margins below the
capped dual-index proxy are negligible and sit inside the heuristic's own optimality band, so the
honest reading is a \emph{match}, not a win; and the multi-echelon improvement is computed against a
different baseline and cost convention than the published A3C figures. The high-penalty,
autocorrelated MMPP2$+$ lost-sales instances at the longest lead times ($L=8,10$) remain won by
classical heuristics, indicating a gap for stationary compact policies under regime-switching demand.
CMA-ES, while far more sample-efficient than A3C here, is still data-hungry relative to data-driven
settings with scarce historical demand, and its population-based evaluation cost grows with the
covariance dimension, so very large parameterizations may need restricted or separable covariance
structures. Promising directions include richer state-dependent decoders for correlated demand,
warm-starting from heuristics in problems beyond dual sourcing, and combining the action-geometry
idea with sample-efficient learners for limited-data regimes.

\appendix

\section{Published vanilla lost-sales validation}

Before using the vanilla lost-sales simulator and heuristic stack for learned-policy comparisons, we
validate the canonical Poisson instance against the published/reference costs used in the lost-sales
literature \citep{Zipkin2008OldSystems,Xin2021CappedBaseStock}. The validation instance has
$h=1$, $p=4$, $L=4$, and $D_t\sim\mathrm{Poisson}(5)$ (repo anchor
\texttt{lit\_poisson\_p4\_l4}). Table~\ref{tab:vanilla-lost-sales-validation} reports the reference
costs and a current Rust rollout check for the heuristic policies
(\texttt{lost\_sales\_heuristics\_all}, horizon $10^5$, seed $123$). This is an implementation
check, not a learned-policy result.

\begin{table}[!htbp]
\centering
\caption{Published/reference vanilla lost-sales validation example for the canonical
Poisson-demand instance. The reproduced heuristic costs use the same long-horizon rollout protocol as
Table~\ref{tab:results-vanilla-lost-sales}.}
\label{tab:vanilla-lost-sales-validation}
\begin{tabular}{lll}
\toprule
Quantity & Published/reference cost & Current code check \\
\midrule
optimal & 4.73 & 4.73 (catalog reference) \\
myopic-1 & 5.06 & 5.0569 \\
myopic-2 & 4.82 & 4.8208 \\
SVBS & 5.83 & 5.8153 \\
\bottomrule
\end{tabular}
\end{table}
\FloatBarrier

The heuristic costs match the published/reference values at the precision used in the benchmark
tables, anchoring the vanilla lost-sales environment and the myopic/SVBS heuristic implementation.
The MMPP2 rows in Table~\ref{tab:results-vanilla-lost-sales} are mean-preserving repo extensions and
are not claimed as published literature rows.

\section{Published fixed-cost lost-sales validation}

Before relying on the fixed-cost simulator and heuristic stack for learned-policy comparisons, we
validate the implementation against the published example of
\citet{Bijvank2015ParametricPolicies}. Their Table~1 reports a small Poisson-demand instance with
$L=2$, $p=14$, $K=5$, $h=1$, $\mu=5$ (the repo anchor
\texttt{bijvank2015\_table1\_l2\_p14\_k5}, exposed through the
\texttt{invman\_rust} fixed-order-cost reference catalog). The current code check uses the Rust
exact average-cost value-iteration solver and exact evaluators for the published heuristic
parameters.

\begin{table}[!htbp]
\centering
\caption{Published fixed-cost lost-sales validation example from
\citet{Bijvank2015ParametricPolicies} and our reproduced exact average costs for the reported
heuristic parameters.}
\label{tab:fixed-cost-validation}
\begin{tabular}{llll}
\toprule
Policy & Reported parameters & Published cost & Our reproduced cost \\
\midrule
optimal & --- & 11.46 & 11.4631 \\
$(s,S)$ & $(17,23)$ & 11.62 & 11.6181 \\
$(s,nQ)$ & $(17,7)$ & 11.56 & 11.5552 \\
modified $(s,S,q)$ & $(17,23,7)$ & 11.50 & 11.4974 \\
\bottomrule
\end{tabular}
\end{table}
\FloatBarrier

The reproduced exact costs match the published values to the precision of the reported table,
anchoring the fixed-cost environment and heuristic implementation to a published benchmark row.

\section{Environment validation against published costs}

Before using the multi-echelon environment for learned-policy comparison, we validate it against
the published \emph{constant base-stock} costs of \citet{vanroy1997ndp} (the simple problem and the
two complex case studies), evaluated at the published $(y^w,y^r)$ levels under the original
Van-Roy-style cost convention (\texttt{van\_roy\_1997} dynamics), with $100$ replications of horizon
$10^5$, no warm-up discard, and the minimum-spread allocation. Table~\ref{tab:me-validation} shows
that all three reproduce within the $2\%$ simulation-protocol tolerance (maximum absolute gap
$1.31\%$).

\begin{table}[!htbp]
\centering
\caption{Reproduction of the published Van Roy constant base-stock costs.}
\label{tab:me-validation}
\begin{tabular}{lcccc}
\toprule
Instance & Levels $(y^w,y^r)$ & Published & Reproduced & Gap \\
\midrule
Simple problem & $(10,16)$ & $51.7$ & $51.72$ & $+0.03\%$ \\
Case study 1 & $(330,23)$ & $1302$ & $1284.9$ & $-1.31\%$ \\
Case study 2 & $(460,22)$ & $1449$ & $1437.6$ & $-0.79\%$ \\
\bottomrule
\end{tabular}
\end{table}
\FloatBarrier

Two caveats keep this an implementation-fidelity check rather than a full literature verification.
First, exact reproduction requires calibrated demand inputs: the simple problem uses a latent mean
of $6.294$ (the effective mean of Van Roy's $\mathcal{N}(5,8^2)$ after rounding and clipping to
non-negative integers), and case study~2 a latent mean of $1.0$ rather than the stated $0$. Second,
we reproduce only the constant base-stock rows, not the A3C learner. The residual $\sim 1.3\%$ is
attributed to unspecified details of Van Roy's original simulation protocol. The paper-faithful
settings used for policy design (Table~\ref{tab:me-params}) instead use the \texttt{gijs\_2022} cost
convention and the Table-3 demand means.

\section{Dual-sourcing comparator and the bounded-DP optimum}
\label{app:ds-validation}

For dual sourcing we benchmark against the capped dual-index heuristic as the optimal proxy rather
than a computed optimum. An exhaustive grid search over the capped-dual-index control
$(s_e,\Delta_r,\bar c_r)$ under the evaluation protocol recovers the published capped-dual-index
parameters and confirms that it is the best static control on every Figure~9 row, consistent with
its published $\le0.11\%$ optimality gap. We do \emph{not} use our own bounded-DP value iteration as
the denominator: at inventory bounds that keep the $l_r=4$ state space tractable, the truncated
bounded-DP cost is mildly inconsistent --- it falls slightly below the simulated heuristic costs ---
and widening the bounds to remove the truncation is computationally prohibitive at $l_r=4$. The
capped dual-index proxy is near-exact to within its published gap and is evaluated on the identical
common-random-number rollout as the learned policies.

% Bibliography embedded inline (natbib-compatible \bibitem labels) so the manuscript is
% self-contained and renders with a single pdflatex pass -- no bibtex step required. The
% same entries are also kept in references.bib for tooling that prefers a .bib database.
\begin{thebibliography}{17}

\bibitem[Bijvank et al.(2015)]{Bijvank2015ParametricPolicies}
M.~Bijvank, S.~Bhulai, and W.~T. Huh.
Parametric replenishment policies for inventory systems with lost sales and fixed order cost.
\emph{European Journal of Operational Research}, 241(2):381--390, 2015.

\bibitem[Bijvank and Vis(2011)]{Bijvank2011LostSalesReview}
M.~Bijvank and I.~F.~A. Vis.
Lost-sales inventory theory: A review.
\emph{European Journal of Operational Research}, 215(1):1--13, 2011.

\bibitem[Boute et al.(2022)]{boute2022deep}
R.~N. Boute, J.~Gijsbrechts, W.~van Jaarsveld, and N.~Vanvuchelen.
Deep reinforcement learning for inventory control: A roadmap.
\emph{European Journal of Operational Research}, 298(2):401--412, 2022.

\bibitem[Hansen and Ostermeier(2001)]{Hansen2001CompletelyStrategies}
N.~Hansen and A.~Ostermeier.
Completely derandomized self-adaptation in evolution strategies.
\emph{Evolutionary Computation}, 9(2):159--195, 2001.

\bibitem[Hansen et al.(2019)]{Hansen2019pycma}
N.~Hansen, Y.~Akimoto, and P.~Baudis.
CMA-ES/pycma on GitHub.
Zenodo, DOI:10.5281/zenodo.2559634, 2019.

\bibitem[Klambauer et al.(2017)]{Klambauer2017Self}
G.~Klambauer, T.~Unterthiner, A.~Mayr, and S.~Hochreiter.
Self-normalizing neural networks.
In \emph{Advances in Neural Information Processing Systems}, volume~30, 2017.

\bibitem[Salimans et al.(2017)]{Salimans2017EvolutionLearning}
T.~Salimans, J.~Ho, X.~Chen, S.~Sidor, and I.~Sutskever.
Evolution strategies as a scalable alternative to reinforcement learning.
\emph{arXiv preprint arXiv:1703.03864}, 2017.

\bibitem[Zipkin(2008)]{Zipkin2008OldSystems}
P.~Zipkin.
Old and new methods for lost-sales inventory systems.
\emph{Operations Research}, 56(5):1256--1263, 2008.

\bibitem[Xin(2021)]{Xin2021CappedBaseStock}
L.~Xin.
Technical Note---Understanding the performance of capped base-stock policies in lost-sales inventory models.
\emph{Operations Research}, 69(1):61--70, 2021.

\bibitem[Van Roy et al.(1997)]{vanroy1997ndp}
B.~Van Roy, D.~P. Bertsekas, Y.~Lee, and J.~N. Tsitsiklis.
A neuro-dynamic programming approach to retailer inventory management.
In \emph{Proceedings of the 36th IEEE Conference on Decision and Control}, volume~4, pages 4052--4057, 1997.

\bibitem[Gijsbrechts et al.(2022)]{gijsbrechts2022drl}
J.~Gijsbrechts, R.~N. Boute, J.~A. Van Mieghem, and D.~J. Zhang.
Can deep reinforcement learning improve inventory management? Performance on lost sales, dual-sourcing, and multi-echelon problems.
\emph{Manufacturing \& Service Operations Management}, 24(3):1349--1368, 2022.

\bibitem[Nahmias and Smith(1994)]{nahmias1994partial}
S.~Nahmias and S.~A. Smith.
Optimizing inventory levels in a two-echelon retailer system with partial lost sales.
\emph{Management Science}, 40(5):582--596, 1994.

\bibitem[Federgruen and Zipkin(1984)]{federgruen1984multilocation}
A.~Federgruen and P.~Zipkin.
Approximations of dynamic, multilocation production and inventory problems.
\emph{Management Science}, 30(1):69--84, 1984.

\bibitem[de Kok et al.(2018)]{dekok2018typology}
T.~de Kok, C.~Grob, M.~Laumanns, S.~Minner, J.~Rambau, and K.~Schade.
A typology and literature review on stochastic multi-echelon inventory models.
\emph{European Journal of Operational Research}, 269(3):955--983, 2018.

\bibitem[Kaynov et al.(2024)]{kaynov2024owmr}
I.~Kaynov, M.~van Knippenberg, V.~Menkovski, A.~van Breemen, and W.~van Jaarsveld.
Deep reinforcement learning for one-warehouse multi-retailer inventory management.
\emph{International Journal of Production Economics}, 267:109088, 2024.

\bibitem[Veeraraghavan and Scheller-Wolf(2008)]{veeraraghavan2008dualindex}
S.~Veeraraghavan and A.~Scheller-Wolf.
Now or later: A simple policy for effective dual sourcing in capacitated systems.
\emph{Operations Research}, 56(4):850--864, 2008.

\bibitem[Sheopuri et al.(2010)]{sheopuri2010dualsourcing}
A.~Sheopuri, G.~Janakiraman, and S.~Seshadri.
New policies for the stochastic inventory control problem with two supply sources.
\emph{Operations Research}, 58(3):734--745, 2010.

\end{thebibliography}

\end{document}
